\documentclass[11pt]{article}
\usepackage[a4paper,margin=23mm]{geometry}
\usepackage[T1]{fontenc}
\usepackage{lmodern,amsmath,amssymb,amsthm,mathtools,microtype,booktabs}
\usepackage[colorlinks=true,linkcolor=blue,citecolor=blue,urlcolor=blue]{hyperref}
\newcommand{\E}{\mathbb E}
\newcommand{\tr}{\operatorname{tr}}
\newcommand{\Wick}{\mathcal W}
\newcommand{\Ncal}{\mathcal N}
\newcommand{\norm}[1]{\left\lVert#1\right\rVert}
\newcommand{\Hh}{\mathcal H}
\newtheorem{theorem}{Theorem}
\newtheorem{lemma}{Lemma}
\title{Fixed-charge gradients and the ordinary charged tail\\
\large CAR channel identities, marked Wick estimates and finite-time summation}
\author{Alan Day \and Chengheng\thanks{AI research assistant (OpenAI).}}
\date{8 November 2026\quad Local research draft}
\begin{document}\maketitle
\begin{abstract}
We replace the full-Fock Riesz step in the ordinary charged-tail argument
by estimates on adjacent charge sectors. Normalized up/down channels
contract the principal row and column gradient terms; their remainders
accumulate linearly along a channel path. A graded derivative shifts
each body/size deficit by one half. This gives a direct marked-Wick
gradient estimate with a square-root-p cost, uniform in polynomial
degree. A rectangular completely positive map estimate controls
products at fixed Schatten orders. The actual covariance residual and
forced Jacobi recurrence then yield a finite-N, degree-uniform ordinary
tail in a central charge window. Combined with the previous fixed-degree
half-filling bridge, this permits finite-time summation in the
previously controlled polynomial preparations.
\end{abstract}

\section{Model, norms and statement}
Use standard CAR, $\{c_i,c_j^\dagger\}=\delta_{ij}$, and
$\Hh_m=\wedge^m\mathbb C^N$, $d_m=\binom Nm$. The Hermitian Gaussian
overlapping p-body Hamiltonian is the same common-disorder model as
in the preceding drafts. Its directed coefficients are
$A_{IJ}=\sqrt{w_N}C_I^\dagger C_J$, $|I|=|J|=p$, with covariance
\[
 C_m I=\E H_m^2,\qquad C_m=w_N\binom mp\binom{N-m+p}p.
\]
Take even N, $\lambda_N=p^2/N\in[\lambda_-,\lambda_+]\subset(0,\infty)$,
and fixed energy normalization $J_s>0$. Along the double-scaling
sequence p tends to infinity. Covariances are bounded above by $C_*$
on all sectors and bounded below on $|m-N/2|\le L_0p$ for each fixed
$L_0$, by the binomial estimates in \cite{operator}.
All norms include Gaussian expectation:
\[
 \norm{X}_{s,m}=(\E\tr_m|X|^s)^{1/s},\qquad
 \tr_m=d_m^{-1}\operatorname{Tr}_{\Hh_m}.
\]
On neutral block families put
\[
 TX=\sum_i c_i^\dagger Xc_i,\quad SX=\sum_i c_iXc_i^\dagger,
 \quad D=Q-T,\quad \bar D=N-Q-S,\quad L=(D+\bar D)/2.
\]
Loss acts on the full family before restriction. $\Ncal$ is the
coupling Gaussian number, not body degree; $Z_n=\Wick_n(H^n)$ is top
coupling chaos. Fixed-charge ensembles have the established background
of \cite{bnr}; the present finite-N bounds concern this specific model.

\begin{theorem}\label{thm:tail}
For fixed finite s and $L_0$, constants $K_s,L_s,b_s$ independent of
N and n satisfy
\begin{align}
 G_{n,N;m}&=C_m^{-n/2}
 [p\Ncal-D-\lambda_N\Ncal(\Ncal-1)]H^n\big|_{\Hh_m},\nonumber\\
 \sup_{|m-N/2|\le L_0p}\norm{G_{n,N;m}}_{s,m}
 &\le K_s(n+1)^{b_s}L_s^n\sqrt{n!}.\label{eq:main}
\end{align}
$C_m$ here is a fixed target scalar applied after loss. For finite S,
the actual series $G_N(t)=\sum_{n\ge0}(it)^nG_{n,N;m}/n!$ has tail
\[
 \sup_N\sup_{|t|\le S}
 \norm{G_N(t)-\sum_{n\le K}(it)^nG_{n,N;m}/n!}_{s,m}
 \le K_s\sum_{n>K}\frac{(n+1)^{b_s}(L_sS)^n}{\sqrt{n!}}
 \longrightarrow0. \tag{T}
\]
Time uses dimensionless $h=H_m/\sqrt{C_m}$ units and $\hbar=1$.
Finitely many initial N are absorbed in the constants.
\end{theorem}

The fixed-degree half-filling bridge is in \cite{central}. The
full-Fock Clifford estimate of \cite{belp}, Theorem 6.8, is not a
fixed-charge theorem; we do not restrict it to a sector. We use
the earlier pure-tag and covariance identities
\cite{diagram,operator,charged}, with their sectorwise bounds,
and supply the missing gradient and product estimates.

\section{Adjacent-sector gradients and channel transport}
Put $Y_{i,m}(X)=[c_i,X]_m:\Hh_m\to\Hh_{m-1}$ and
\[
 g_s^C(X;m)=\norm{(\sum_iY_{i,m}^\dagger Y_{i,m})^{1/2}}_{s,m},
 \qquad
 g_s^R(X;m)=\norm{(\sum_iY_{i,m}Y_{i,m}^\dagger)^{1/2}}_{s,m-1}.
\]
Let $g_s$ be the maximum of both quantities for X and $X^\dagger$.
Raising derivatives are covered by adjoints. Interval suprema
include the neighboring sectors of their edges. We work at $s\ge2$;
higher even orders imply other fixed orders by monotonicity.

The channels $(tX)_k=(TX)_k/k$, $(sX)_k=(SX)_k/(N-k)$ are unital
and preserve normalized sector trace, by cyclicity and
$d_{k-1}/d_k=k/(N-k+1)$. Thus they contract every Schatten order,
also after Gaussian averaging. The channel s and norm exponent s
are distinguished by context.

\begin{lemma}\label{lem:channel}
On an interior edge $m\to m-1$, with $d=N-m+1$,
\begin{align}
 [c_i,tX]_m&=-\frac1m\sum_jc_j^\dagger Y_{i,m-1}(X)c_j
       +\frac1m(X-tX)_{m-1}c_i,\label{eq:tgrad}\\
 [c_i,sX]_m&=-\frac1d\sum_jc_jY_{i,m+1}(X)c_j^\dagger
       +\frac1d c_i(sX-X)_m.\label{eq:sgrad}
\end{align}
The principal terms contract both row and column square-function
norms with factors $\sqrt{(m-1)/m}$ and $\sqrt{(d-1)/d}$.
\end{lemma}
\begin{proof}
CAR gives $[c_i,TX]=Xc_i-\sum_jc_j^\dagger[c_i,X]c_j$ and
$[c_i,SX]=-c_iX-\sum_jc_j[c_i,X]c_j^\dagger$.
The left/right normalized denominators are m and $m-1$ for T,
and $d-1$ and d for S. Keeping them distinct gives the identities.

For arbitrary $Y_i:\Hh_{m-1}\to\Hh_{m-2}$, let
$P_i=\sum_jc_j^\dagger Y_i c_j$,
$Q_C=\sum_iY_i^\dagger Y_i$, $Q_R=\sum_iY_iY_i^\dagger$.
Row/column operator Cauchy--Schwarz gives
\[
 \sum_iP_i^\dagger P_i\le(m-1)TQ_C,\qquad
 \sum_iP_iP_i^\dagger\le mTQ_R.
\]
The first T targets m and the second targets $m-1$. After division
by $m^2$, both yield $(m-1)/m$ times normalized t. For
$Y_i:\Hh_{m+1}\to\Hh_m$ and $P_i=\sum_jc_jY_i c_j^\dagger$,
\[
 \sum_iP_i^\dagger P_i\le dSQ_C,\qquad
 \sum_iP_iP_i^\dagger\le(d-1)SQ_R.
\]
After division by $d^2$, both yield $(d-1)/d$ times normalized s.
Positivity, Schatten monotonicity, and channel contraction at order
$s/2$ prove the bounds. The proof also applies to $X^\dagger$.
\end{proof}

For the remainder $U_{m-1}c_i/m$, the column and row norms are,
respectively, bounded by
\[
 m^{-1/2}\norm U_{s,m-1},\qquad
 \sqrt{N-m+1}\norm U_{s,m-1}/m,
\]
with equality in the row expression. For $c_iV_m/d$ they are
$\sqrt m\norm V_{s,m}/d$ and at most $\norm V_{s,m}/\sqrt d$.
These follow by summing the Gram matrices before taking norms.
On $[N/4,3N/4]$, each step therefore adds at most
$8N^{-1/2}\sup\norm X_s$, since $U=X-tX$, $V=sX-X$.
For a normalized path P of length $\ell$, all of whose required
edges and sectors stay in that interval, induction gives
\[
 g_s(PX;m)\le\sup_{\rm reachable}g_s(X)
   +\frac{8\ell}{\sqrt N}\sup_{\rm reachable}\norm X_s. \tag{P}
\]
Intermediate norms contract. No factor greater than one is raised
to $\ell$. For a scalar charge function f the extra derivative is
$(f_m-f_{m-1})X_{m-1}c_i$, so
\[
 g_s(fX;m)\le\sup|f|\,g_s(X;m)
    +\sqrt N\sup|\Delta f|\sup\norm X_s. \tag{F}
\]

\section{The graded shift and the pure marked gradient}
For homogeneous parity $a\in\{0,1\}$ define algebraically
\[
 D_gX=\tfrac12\{Q,X\}-(-1)^a\sum_i c_i^\dagger Xc_i,\quad
 \partial_bX=bX-(-1)^aXb,\quad
 L_g=D_g-\tfrac12\sum_i\partial_{c_i^\dagger}\partial_{c_i}.
\]
On neutral even fields these agree with D,L. No odd-field CP
property is assumed. A normal-ordered monomial of total degree d
is a $D_g$ eigenvector with eigenvalue $d/2$. A graded derivative
deletes one matching letter with its position sign. Such monomials
form a basis; hence for $b=c_i$ or $c_i^\dagger$,
\[
 [\partial_b,D_g]=\tfrac12\partial_b,\qquad
 [\partial_b,L_g]=\tfrac12\partial_b. \tag{G}
\]
For L, use the anticommutation of graded derivatives: their
quadratic contraction commutes with each derivative. For an ordered
list $T_j\in\{D,L\}$,
\[
 \partial_b\prod_{j=1}^k(np-T_j)Z_n
 =\prod_{j=1}^k(np-\tfrac12-T_{j,g})\,\partial_bZ_n. \tag{S}
\]
The order is preserved; commutativity of deficits is not required.

\begin{lemma}\label{lem:pure}
For fixed k,s and a fixed interior charge interval, uniformly in N,n,
\[
 g_s\!\left(\prod_{j=1}^k(np-T_j)Z_n;m\right)
 \le K_{k,s}\sqrt p\,(n+1)^{2k+1}L_{k,s}^n\sqrt{n!}. \tag{W}
\]
The field without the gradient has the earlier bound
$K_{k,s}(n+1)^{2k}L_{k,s}^n\sqrt{n!}$.
\end{lemma}
\begin{proof}
We extend the restricted-covariance proof of \cite{diagram}.
The graded Leibniz rule for $\partial_bZ_n$ chooses one of n
original coefficient slots and deletes a matching letter, imposing
one membership in its original I or J label. The remaining total
degree is $2np-1$, so (S) uses exactly the appropriate deficit.
The graded product rule is
\[
 D_g(XY)=(D_gX)Y+X(D_gY)
       +(-1)^{|X|}\sum_i(\partial_{c_i^\dagger}X)(\partial_{c_i}Y).
\]
Each subsequent deficit either multiplies a term by its number
of already deleted pairs or creates a new two-letter contraction;
L also permits an own-slot I/J contraction. After k operations
there are at most $K_k(n+1)^{2k+1}$ terms, with coefficients
depending only on k. Deleting an already deleted letter is zero.
All surviving memberships remain requirements on the original
sets; (S) avoids differentiating inserted sandwich letters.

Expand a column square-function moment
$\E\tr_m(\sum_iY_i^\dagger Y_i)^r$ (order $2r$).
There are $2r$ differentiated Wick blocks, at most $2rk$ pure
tags, and r summed gradient indices, each occurring twice.
Fix the Gaussian pairing, slot choices, and partition of all
site variables into equal sites. Let A count distinct sites
carrying a pure tag and B sites carrying gradients only.
Every tagged site requires two distinct underlying p-set labels.
For a cross-slot tag its two slots are in one Wick block and
cannot be paired together; an own-slot tag uses independent I,J.
Identifications elsewhere cannot remove those two requirements.
Every gradient-only site needs at least one membership, and $B\le r$.

For fixed membership sets $T_I,T_J$, the occupation-basis column
covariance is bounded by
\[
 C_*\frac{(p)_{|T_J|}}{(m)_{|T_J|}}
       \frac{(p)_{|T_I|}}{(N-m+p)_{|T_I|}}
 \le C_*(K_{k,r}p/N)^{|T_I|+|T_J|}. \tag{R}
\]
If an I site is occupied but not removed by J the contribution
is zero, so dropping that extra restriction is an upper bound.
The row bound swaps I,J. Combine conditions at both ends of a
Gaussian chord before applying (R). Derived coefficients are
finite sums of fixed CAR sandwiches of the restricted original
coefficient; their number depends only on k,r and all insertions
contract.

Scan a chord word with column creation and row annihilation maps
for Gaussian-index ancillas. Crossings only permute open ancillas
unitarily. The contraction norm is bounded by the product of
endpoint covariance square roots, keeping CAR signs throughout.
Each independent set/site condition is counted once per chord.
Summing site assignments therefore costs at most
\[
 N^{A+B}(p/N)^{2A+B}=\lambda_N^A p^B\le K_{k,r}p^r. \tag{C}
\]
Additional memberships improve the estimate. The numbers of site
partitions and intermediate charge shifts depend only on k,r,
not n. For sufficiently large N all intermediate charges remain
in a wider fixed interior interval. Small N have bounded p and
CAR norms and are absorbed with polynomial degree costs.

There are at most $(2rn-1)!!$ Gaussian pairings, with
$[(2rn-1)!!]^{1/(2r)}\le(2r)^{n/2}\sqrt{n!}$.
Taking the root gives the stated slot factor and $\sqrt p$ from
(C). The row moment has the same argument on sector $m-1$;
adjoints and reversed orientations give the same estimates.
Higher even orders imply all fixed orders. Omitting the gradient
gives $B=0$ and the previous pure-tag estimate.
\end{proof}

\section{Products without an ancilla dimension loss}
For $\Phi(X)=\sum_iV_iXV_i^\dagger$ from an a-dimensional sector
to a b-dimensional sector, the normalized-trace adjoint is
$\Phi^\sharp(Y)=(a/b)\sum_iV_i^\dagger YV_i$.
For positive X,Y, spectral decomposition and scalar Holder with
weights $\sum_i|(V_i)_{uv}|^2/b$ imply
\[
 \E\tr_bY\Phi(X)\le
 (\E\tr_aX^u\Phi^\sharp(I))^{1/u}
 (\E\tr_bY^v\Phi(I))^{1/v},\qquad u^{-1}+v^{-1}=1.
\]
The map and inputs may share the disorder. With $u=2r$, duality
at order r and a second trace Holder give
\[
 \norm{\Phi(X)}_{r,b}\le
 4\norm X_{4r,a}\norm{\Phi^\sharp(I)}_{2,a}^{1/(2r)}
                   \norm{\Phi(I)}_{2r-1,b}^{1-1/(2r)}. \tag{CP}
\]
The factor four allows arbitrary X by splitting its real and
imaginary positive/negative parts; positive X needs no factor.
At $r=1$ the dual variable is bounded and the same formula holds.

Apply (CP) to the arrays $Y_i(X)$ and their adjoints. Adjacent
dimension ratios are bounded in the interior. The ordinary
Leibniz rule, with (CP) for inside Gram multipliers and rectangular
Holder for outside multipliers, gives
\begin{align*}
 g_s(XY;m)\le K_s\bigl[
 &g_{4s}(X;m)\sup_{k=m-1,m}\norm{Y_k}_{4s,k}\\
 &+\sup_{k=m-1,m}\norm{X_k}_{4s,k}g_{4s}(Y;m)\bigr]. \tag{B}
\end{align*}
For example the row Gram of $Y_i(X)Y_m$ is
$\Phi_X(Y_mY_m^\dagger)$; its column Gram is
$Y_m^\dagger(\sum_iY_i(X)^\dagger Y_i(X))Y_m$.
Use (CP) at $s/2$ for the former and Holder for the latter.
The other derivative term and adjoints cover all cases. Thus
no trace dimension of $I_N\otimes Y$ has been discarded.

For the body-loss cross term,
\[
 B_D(X,Y)=\sum_i[c_i^\dagger,X][Y,c_i],\qquad
 D(XY)=(DX)Y+X(DY)-B_D(X,Y),
\]
rectangular Holder yields
$\norm{B_D(X,Y)}_{r,m}\le g_{2r}(X;m)g_{2r}(Y;m)$.
The reverse loss uses row norms on edge $m+1$, and L is their
average. No general fixed-charge Riesz theorem is used.

\section{Residuals, ordinary powers and the forced recurrence}
Use nested intervals
$[2N/5,3N/5]\subset[N/3,2N/3]\subset[N/4,3N/4]$.
For large N, $p\le N/100$, leaving room for paths and a fixed
number of neighboring sectors. Write
$F_s(n)=K_s(n+1)^bL_s^n\sqrt{n!}$, allowing constants and b to
increase a fixed number of times.

The covariance and scalar approximation are
\begin{align*}
 \Phi_m(X)&=w_N\sum_{I,J}C_I^\dagger C_JX C_J^\dagger C_I,
 &\Psi_m&=\Phi_m/C_m,\\
 \theta_{v,m}&=\prod_{j=0}^{p-1}
 (1-vp/(m-j))_+(1-vp/(N-m+p-j))_+,
 &A_v&=p(\Phi-C\theta_v)Z_v.
\end{align*}
$\Psi$ is the exact normalized path of p down and p up channels
\cite{operator}. If $vp\le N/8$, its telescoping identity is
\[
 A_v=pC\sum_{h=1}^{2p}f_h(Q)P_hE_h,\qquad
 E_h\in\{(vp-D)Z_v,(vp-\bar D)Z_v\}. \tag{A}
\]
Each $P_h$ has length at most $2p$; $f_h$ is one inverse charge
denominator times remaining scalar factors, with exact shifts.
The sectorwise bounds in \cite{charged} are
\[
 |f_h|\le K/N,\quad |\Delta f_h|\le K(1+v\lambda_N)/N^2,
 \quad C\le C_*,\quad |\Delta C|\le Kp/N.
\]
Lemma \ref{lem:pure}, (P) and (F) give
\[
 g_s(A_v)/\sqrt p\le F_s(v). \tag{A1}
\]
Indeed the principal coefficient sum is $O(p^2/N)$ and the
path cost is $O(p/\sqrt N)$ before division by $\sqrt p$.
After that division the scalar-gradient contribution is at most
$Kp^{3/2}[p+1+v\lambda_N]/N^{3/2}$ times the pure envelope,
with only polynomial v growth. If $vp>N/8$, use the coarse
$g_s(X)\le K\sqrt N\sup\norm X_s$ and
$\norm{A_v}_s\le Kp\norm{Z_v}_s$. Since
$p\le8\lambda_Nv$, $N\le64\lambda_Nv^2$, its cost divided by
$\sqrt p$ is also polynomial in v.

The previous sectorwise derivative calculation gives
\[
 \norm{A_v}_s+\norm{DA_v/p}_s+\norm{LA_v/p}_s\le F_s(v). \tag{A2}
\]
Its inputs, which did not use Riesz, are the exact relations
\[
 [D,t]=t-t^2,\quad
 [D,s]=-s+\frac Q{N-Q}ts-\frac{2Q-N}{N-Q}I,\quad
 [D,f]=(f(Q)-f(Q-1))T
\]
and reverse-loss versions. Expanding a path gives
$O(\ell/p)$ for $[D/p,P]$ and $[L/p,P]$. Repeated pure
deficits control derivatives of $E_h$ in (A); the scalar
differences control the remaining terms. The large-v case uses
the same $N=O(v^2)$ estimate.

All-sector ordinary Gaussian moments obey
$\norm{H^w}_s\le L_s^w\sqrt{w!}$. Lemma \ref{lem:pure} at
$n=1,k=0$ gives both force Grams of $[c_i,H]$ of order p.
Hence (CP) bounds each forward/reverse force map by
$K_rp\norm X_{4r}$. Expanding
\[
 [c_i,H^w]=\sum_{a+b=w-1}H_{m-1}^a[c_i,H]_mH_m^b
\]
and applying this force estimate to inside Gram multipliers
gives $g_s(H^w)/\sqrt p\le F_s(w)$ at fixed higher norm orders.
For D the exact expansion is
\[
 DH^w=pwH^w-\sum_{a+b+c=w-2}H^a\Gamma_H(H^b)H^c,\qquad
 \Gamma_H(X)=\sum_i[c_i^\dagger,H]X[H,c_i].
\]
For L use single-slot LH and averaged forward/reverse cross
terms. Pure $n=1$ controls $LH/p$. Fixed-order Holder and (CP)
then give
\[
 \norm{H^w}_s+\norm{DH^w/p}_s+\norm{LH^w/p}_s
              +g_s(H^w)/\sqrt p\le F_s(w). \tag{H}
\]
Only $O(w^2)$ terms occur, each with at most three powers and
one force map. Norm orders are not doubled w times.

Define the finite-N Jacobi polynomials and actual error by
\begin{gather*}
 Q_{j+1,m}(x)=xQ_{j,m}(x)-b_{j,m}Q_{j-1,m}(x),
 \qquad b_{j,m}=C_m\sum_{v<j}\theta_{v,m},\\
 \eta_j=p[Z_j-Q_j(H)],\qquad Q_0=1,\qquad Q_1=x.
\end{gather*}
The exact forced recurrence \cite{charged} is
\[
 p r_\ell=\sum_{v+w=\ell-1}\Wick_{v+w}(A_vH^w),\qquad
 \eta_j=-\sum_{\ell<j}P_{j,\ell,m}(H)\,p r_\ell. \tag{J}
\]
$P_{j,\ell}$ is the scalar Jacobi propagator, of degree at most
$j-\ell-1$. Its coefficient sums grow at most exponentially,
since $b_{j,m}\le B$ on interior intervals. Their charge
differences satisfy
$N\sum|\Delta f|/p\le\operatorname{poly}(j)K^j$, from
$|\Delta b_{j,m}|\le K[p+j^2]/N$ and path product differences.

The OU/Chebyshev estimate
$\norm{\Wick_dX}_s\le2^d\norm X_s$ for degree-d Gaussian
polynomials \cite{operator} applies unchanged to row/column
arrays: conditional Jensen uses their Banach norms with the
specified sector trace. Physical derivatives commute with
coupling OU and top-chaos projection.

Apply (A1), (A2), (H), (B) and the loss cross bound first to
$A_vH^w$, then top project and sum. Fully expand $P_{j,\ell}(H)$
and apply the same estimates to $H^a(p r_\ell)$, retaining
coefficient differences. There are two product layers,
independent of j. Norm orders remain fixed multiples of s.
Use $\sqrt{a!v!w!}\le\sqrt{j!}$, exponential coefficient
sums and polynomially many degree choices to obtain
\[
 \norm{\eta_j}_s+\norm{D\eta_j/p}_s+\norm{L\eta_j/p}_s
               +g_s(\eta_j)/\sqrt p\le F_s(j). \tag{E}
\]
For coefficient gradients (F) uses
$\sqrt N|\Delta f|/\sqrt p\le\sqrt{p/N}\operatorname{poly}(j)K^j$;
for D/L it uses $N|\Delta f|/p$. All needed neighbors fit the
nested intervals; no sharp charge cutoff is differentiated.

\section{Ordinary powers and selected time fields}
Write $H^n=\sum_jh_{nj,m}Q_{j,m}(H)$ using nonnegative weighted
Jacobi paths. The exact covariance ratio and the Lipschitz
bound for the positive part in $\theta_v$ give \cite{charged}
\[
 \sum_j|h_{nj,m}|\le K^n,\qquad
 \sum_jN|\Delta h_{nj,m}|
 \le K(n+1)^3L^n(1+|X_N|),\quad X_N=(m-N/2)/p. \tag{I}
\]
Explicitly,
\[
 \frac{C_m}{C_{m-1}}-1
 =\frac{p(N-2m+p+1)}{(m-p)(N-m+p+1)},\qquad
 N|\Delta\theta_{v,m}|\le K\lambda_Nv.
\]
The second bound includes positive-part cutoffs. The factor
$1+|X_N|$ is bounded on our central window; it cannot be
discarded on a macroscopic charge interval.

Keeping every charge difference yields
\begin{align}
 (p\Ncal-D)H^n=\sum_j\bigl[&
 h_{nj}(pj-D)Z_j-(\Delta h_{nj})TZ_j-h_{nj}\Ncal\eta_j
 \nonumber\\
 &+h_{nj}D\eta_j/p+(\Delta h_{nj})T\eta_j/p\bigr]. \tag{O}
\end{align}
Pure deficit bounds, (I) and normalized T control the first two
terms. The third uses the degree-j Gaussian number multiplier
and (E), with Chebyshev cost at most $jK^j$. The last two use
(E), (I) and adjacent-sector contraction. The counterterm
$\lambda_N\Ncal(\Ncal-1)H^n$ has the same multiplier estimate
with an extra polynomial in n. Thus \eqref{eq:main} holds before
the target factor $C_m^{-n/2}$; its lower bound changes only
$L_s$. Finitely many small N have bounded p and physical
losses, hence the same Gaussian factorial envelope. Summing
proves (T), also for fixed time derivatives and fixed exterior
energy powers. This completes Theorem \ref{thm:tail}.

At $m=N/2$, combine (T) with the same-copy fixed-degree bridge
\cite{central}. Exact energy pinching contracts every fixed
sector Schatten norm. Each previously controlled fixed
polynomial preparation P has a finite-N orbit density satisfying
$\sup_N\norm{R_{P,N}}_a<\infty$ at each fixed a \cite{prep}.
Holder transfers the actual tail to that preparation, with the
degree of P fixed. Fixed products in a covariant block use
a fixed number of Holder factors. Take N to infinity at fixed
truncation, then remove the truncation with (T). This gives the
complete ordinary charged time series and its exactly pinched
version in that preparation class.

The complete selected bilocal still requires a sectorwise audit
of the actual neutral/rank filter and memory-return estimates.
Spatial/mass generators and finite-G dynamics remain open.
The earlier full-trace bilocal construction is neither retracted
nor promoted to a selected-state result by this theorem.

\section{Finite checks and provenance}
The saved channel check uses $N=3,4,5,6$, three non-selfadjoint
integer neutral fixtures per N and their transposes. It verifies
1,032 gradient and 216 unital/normalized-trace identities exactly
after clearing denominators. A further 168 PSD and 324 Schatten
checks are floating-point diagnostics, not proofs of the lemma
or its path corollary.

The new graded check exhausts 340 matrix units and 340 normal
monomials at $N=1,2,3,4$. It verifies 5,008 derivative
commutators, 840 repeated-deficit shifts, 340 Euler identities
and 2,504 signed deletions, 1,252 of them nonzero. Doubled maps
use exact integer arithmetic. These fixtures check algebra and
signs, not degree-uniform Gaussian convergence. The general
estimates rely on the contraction and norm arguments above,
not numerical extrapolation. Successful scientific outputs
are preserved without rerunning their drivers.

\begin{thebibliography}{9}
\bibitem{belp} L. Ben Efraim and F. Lust-Piquard,
\emph{Poincare type inequalities on the discrete cube and in the CAR algebra},
\href{https://arxiv.org/abs/math/0702233}{arXiv:math/0702233v1},
Section 6.4, Theorem 6.8.
\bibitem{bnr} M. Berkooz, V. Narovlansky and H. Raj,
\emph{Complex Sachdev--Ye--Kitaev model in the double scaling limit},
\href{https://arxiv.org/abs/2006.13983}{arXiv:2006.13983v2}, Section 5.
\bibitem{diagram} Alan Day and Chengheng, \emph{Pure Wick tagged
contractions}, local proof, 7 November 2026, Sections 1--4.1;
\texttt{susskind-bilocal-fourth-tail-diagram-draft.md}.
\bibitem{operator} Alan Day and Chengheng, \emph{OU projections
and covariance CP paths}, local proof, 7 November 2026, Sections 1--5;
\texttt{susskind-bilocal-fourth-tail-operator-draft.md}.
\bibitem{charged} Alan Day and Chengheng, \emph{Ordinary charged
tail on the full Fock trace}, local proof, 7 November 2026,
Sections 2--3 and 5--7;
\texttt{susskind-bilocal-fourth-tail-charged-proof.md}.
\bibitem{central} Alan Day and Chengheng, \emph{Ordinary charged
fields at fixed half filling}, local draft, 8 November 2026.
\bibitem{prep} Alan Day and Chengheng, \emph{Fixed polynomial
Slater preparations}, local draft, 8 November 2026.
\end{thebibliography}
\end{document}
